% Notas de aula — Capítulo 22: Óptica Atmosférica
% Curso: Meteorologia Prática (baseado em R. Stull, Practical Meteorology, CC BY-NC-SA 4.0)
% Caixas verdes = conteúdo literal do quadro; linhas verdes = encenação.
\documentclass[11pt]{article}
\usepackage{../comum/preambulo}
\graphicspath{{../figuras/cap22/}}

\title{Capítulo 22 --- Óptica Atmosférica\\
{\large Notas de aula (roteiro de quadro-negro)}}
\author{Curso de Meteorologia Prática}
\date{}

\begin{document}
\maketitle
\thispagestyle{fancy}

\noindent\textbf{Plano sugerido:} 2 aulas de 2\,h.
\emph{Aula 1:} \S\ref{sec:arco}--\S\ref{sec:halo} (arco-íris e halos —
gotas e cristais como instrumentos ópticos).
\emph{Aula 2:} \S\ref{sec:ceu}--\S\ref{sec:miragem} (espalhamento, as
cores do céu, refração e miragens).

\section{O capítulo-sobremesa}\label{sec:intro}

\begin{noquadro}[abertura]
\textbf{Cap.~22 --- Óptica Atmosférica}\\[2pt]
o céu é um laboratório de óptica: cada fenômeno luminoso é um
\emph{diagnóstico}\\
arco-íris $=$ gotas grandes \quad halo $=$ cristais de gelo (cirrus!)
\quad miragem $=$ perfil de $T$\\[2pt]
ferramentas: Snell, dispersão, espalhamento, e a velha
Clausius--Clapeyron das gotas do Cap.~7
\end{noquadro}

Fechamento perfeito do curso para físicos: só óptica de graduação
aplicada às partículas que os Caps.~4--7 construíram. E cada fenômeno
\emph{informa}: halo anuncia cirrus da frente quente que chega
\javimos{12}{nuvens frontais adiantadas}; o tamanho das gotas decide
se há arco ou glória \javimos{7}{distribuição de tamanhos}.

\section{Arco-íris}\label{sec:arco}

\quadro{Desenhar a gota com o raio: refração, UMA reflexão interna,
refração — e definir o desvio $D$. Anunciar a pergunta: por que 42°?}

\begin{formadif}[a cáustica dos 42° (T1; Caso~1)]
\[ D = 180° + 2\theta_i - 4\theta_r,\qquad
\sin\theta_i = n\sin\theta_r. \]
$dD/d\theta_i = 0$ em $\cos^2\theta_i = (n^2-1)/3$: um
\emph{intervalo} de raios sai no mesmo ângulo — a cáustica. Com
$n = 1{,}33$: arco a $42{,}4°$ (vermelho) e $40{,}7°$ (violeta) do
ponto antissolar.
\end{formadif}

\begin{noquadro}[o dicionário do arco]
vermelho FORA, violeta DENTRO (dispersão: $n$ do violeta maior)\\
secundário a $\sim$51°, cores INVERTIDAS (duas reflexões internas)\\
banda escura de Alexandre entre os dois (nenhum raio sai ali)\\
brilho supranumerário: interferência — o arco é também experimento de
ondas
\end{noquadro}

O arco exige gotas grandes ($\gtrsim$0,5\,mm — chuva
\javimos{7}{gotas de chuva}); gotículas de nuvem ($\sim$10\,$\mu$m)
fazem a \emph{glória} e o arco de névoa branco (difração dominando). O
Monte Carlo do Caso~1 mostra a concentração de raios no ângulo do
arco; o N1 mede a largura com $n(\lambda)$ real.

\section{Halos: os cristais como prismas}\label{sec:halo}

\quadro{Desenhar o hexágono de gelo como prisma de 60° e o caminho do
raio pelas faces alternadas.}

\begin{formadif}[o mínimo dos 22° (T2; Caso~2)]
\[ \delta_{min} = 2\arcsin\!\left(n\sin\frac{A}{2}\right) - A
\qquad (A = 60°,\ n = 1{,}31)\ \Rightarrow\ 21{,}8°. \]
Cristais orientados ao acaso acumulam-se no mínimo raso: o anel de
22°, escuro por dentro (nenhum desvio menor existe), borda interna
avermelhada (N2).
\end{formadif}

\begin{noquadro}[a família dos halos]
halo 22°: colunas/placas ao acaso — o mais comum\\
parélios (sun dogs): placas orientadas horizontais, à altura do Sol\\
pilar de luz: reflexo nas bases das placas\\
halo 46° e arcos tangentes: faces de 90° e orientações especiais\\[2pt]
todos $=$ \textbf{cirrus} \javimos{6}{nuvens altas de gelo}: halo hoje,
frente quente amanhã \javimos{12}{sequência Ci$\to$Cs$\to$As}
\end{noquadro}

\section{Por que o céu é azul (e o poente vermelho)}\label{sec:ceu}

\quadro{A pergunta mais antiga da turma finalmente respondida com o
dipolo de Rayleigh — e com o Beer do Cap.~2 fazendo o poente.}

\begin{formadif}[Rayleigh (T3; Caso~3)]
Molécula $\ll \lambda$ como dipolo forçado: potência irradiada
$\propto \omega^4$, logo
\[ \sigma \propto \lambda^{-4}
\quad\Rightarrow\quad
\frac{\sigma_{azul}}{\sigma_{verm}} =
\left(\frac{650}{450}\right)^4 \approx 4{,}4. \]
Céu $=$ luz \emph{espalhada} (azul vence); Sol poente $=$ luz
\emph{transmitida} por massa de ar $1/\sin\Psi$
\javimos{2}{Beer e massa de ar}: o azul foi roubado no caminho.
\end{formadif}

\begin{noquadro}[o dicionário das cores]
céu azul (não violeta): espectro solar $+$ olho $+$ ozônio\\
nuvem branca: gotas $\gg\lambda$ espalham tudo igual (Mie)\\
poente vulcânico roxo: aerossol estratosférico
\javimos{21}{Pinatubo} (N3)\\
céu polarizado a 90° do Sol (T5): o dipolo não irradia no próprio eixo
\end{noquadro}

\section{Refração: miragens, Fata Morgana e o flash verde}
\label{sec:miragem}

\quadro{Escrever a equação do raio e desenhar os dois casos: asfalto
quente (raio curva para cima) e mar frio (para baixo).}

\begin{formadif}[o raio que curva (T4; Caso~4)]
$n$ segue a densidade \javimos{1}{lei dos gases}: ar quente $=$ $n$
menor. Raio quase horizontal:
\[ \frac{d^2z}{dx^2} \simeq \frac{1}{n}\frac{dn}{dz},\qquad
R_{curv} = \left|\frac{n}{dn/dz}\right|. \]
Asfalto ($dT/dz \sim -2$\,K/m no 1º metro): o ``lago'' que é céu
refletido. Inversão marítima ($dT/dz > 0$
\javimos{5}{inversões}): miragem superior, \emph{looming}, Fata
Morgana; $dT/dz \approx +0{,}11$\,K/m guia o raio na curvatura da
Terra (N4).
\end{formadif}

A refração padrão também: achata o Sol poente, adianta o nascer em
$\sim$2\,min, e — com a dispersão — produz o \textbf{flash verde} no
último gomo do disco. Cintilação das estrelas: turbulência do Cap.~18
\javimos{18}{turbulência} vista como ruído de fase.

\section{Síntese da aula --- e do curso}

\begin{noquadro}[mapa final]
\[ \cos^2\theta_i = \frac{n^2-1}{3}\ (42°) \qquad
\delta_{min} = 2\arcsin(n\sin\tfrac{A}{2}) - A\ (22°) \]
\[ \sigma \propto \lambda^{-4} \qquad
\frac{d^2z}{dx^2} = \frac{1}{n}\frac{dn}{dz} \]
olhar o céu e LER: gotas, cristais, perfis de temperatura — o curso
inteiro num pôr do sol
\end{noquadro}

\section{Exercícios complementares}\label{sec:exercicios}

Soluções (professor) em \texttt{cap22\_solucoes.ipynb}.

\subsection*{Teóricos}

\begin{enumerate}[label=\textbf{T\arabic*.},itemsep=4pt]
  \item Minimize o desvio do raio na gota e mostre que
        $\cos^2\theta_i = (n^2-1)/3$; calcule os ângulos do arco para
        o vermelho e o violeta.
  \item Deduza $\delta_{min} = 2\arcsin(n\sin A/2) - A$ pela simetria
        do percurso no prisma e avalie para o gelo (60°).
  \item Obtenha o $\lambda^{-4}$ de Rayleigh pelo argumento do dipolo
        forçado e explique por que o céu não é violeta e a nuvem é
        branca.
  \item Mostre que o raio de curvatura de um raio quase horizontal é
        $R = n/|dn/dz|$ e calcule o gradiente térmico que iguala $R$
        ao raio da Terra.
  \item Explique a polarização do céu a 90° do Sol pelo padrão de
        radiação do dipolo, e por que o arco-íris também é
        polarizado.
\end{enumerate}

\subsection*{Numéricos (no notebook)}

\begin{enumerate}[label=\textbf{N\arabic*.},itemsep=4pt]
  \item Com $n(\lambda)$ da água, calcule a largura do arco primário
        e a posição do secundário.
  \item Monte Carlo do halo em duas cores: separação angular e a
        borda vermelha interna.
  \item Acrescente aerossol vulcânico ao crepúsculo e explique os
        poentes roxos pós-Pinatubo.
  \item Inverta o perfil térmico (mar frio) e produza a miragem
        superior/looming; estime o alcance extra do horizonte.
\end{enumerate}

\vspace{1em}
\noindent\rule{\linewidth}{0.4pt}\\
{\scriptsize Material derivado de R.~Stull, \emph{Practical Meteorology}
(CC BY-NC-SA 4.0); este documento herda a mesma licença.}

\end{document}
